% SEGMENT OLP-0395-S01
% Part: many-valued-logic
% Chapter: three-valued-logics
% Section: kleene

\documentclass[../../../include/open-logic-section]{subfiles}

\begin{document}

\olfileid{mvl}{thr}{skl}

\olsection{க்ளீனி தருக்கங்கள்}

% SOURCE CORRECTION TA-MVL-011: repair the missing “as” and malformed “the if” in the computational motivation.
ஸ்டீவன் க்ளீனி, மெய்மதிப்புகளைக் கணிப்பீட்டுச் செயல்முறைகளின் விளைவுகளாகக்
கருதும் ஒரு தருக்கத்தால் தூண்டப்பட்ட இரு மும்மதிப்புத் தருக்கங்களை
அறிமுகப்படுத்தினார்: ஒரு செயல்முறை $\True$ அல்லது $\False$ ஐ விளைவாகத்
தரலாம்; ஆனால் அது முடிவுறாமலும் போகலாம். அப்போது அதற்குரிய மெய்மதிப்பு
வரையறுக்கப்படாதது; அதை மெய்மதிப்பு~$\Undef$ குறிக்கிறது.

ஒரு கூற்று~$!A$ இன் மறுப்பைக் கணிக்க, முதலில்~$!A$ இன் மதிப்பைக் கணித்து,
பின்னர் அதன் எதிர்மதிப்பைத் திருப்பித் தருவோம். $!A$ இனது கணிப்பீடு
முடிவுறாவிட்டால், முழுச் செயல்முறையும் முடிவுறாது: ஆகவே $\Undef$ இன்
மறுப்பு~$\Undef$.

ஓர் இணையல்~$!A \land !B$ ஐக் கணிப்பதற்கு இரண்டு வாய்ப்புகள் உள்ளன:
முதலில்~$!A$ ஐயும், பின்னர்~$!B$ ஐயும் கணித்து, இரண்டின் விளைவும்~$\True$
எனில் முடிவை~$\True$ என்றும், இல்லையெனில்~$\False$ என்றும் கொள்ளலாம்.
இரண்டு கணிப்பீடுகளில் ஏதேனும் ஒன்று முடிவுறத் தவறினால், முழுச்
செயல்முறையும் முடிவுறத் தவறும். எனவே இந்நேர்வில் ஓர் இணையுறுப்பு
வரையறுக்கப்படாததாக இருந்தால், இணையலும் அவ்வாறே இருக்கும். பிரிப்பிணைவுக்கும்
இதே பொருந்தும்.

ஆனால் $!A$, $!B$ ஆகியவற்றை இணைநிலையில் மதிப்பிட முடியுமானால் இன்னும்
சிறப்பாகச் செய்யலாம். இரு செயல்முறைகளில் ஒன்று முடிவுற்று~$\False$ ஐத்
திருப்பித் தந்தால், விடை பொய்யாகவே இருக்க வேண்டும் என்பதால் நிறுத்தலாம்.
ஆகவே இந்நேர்வில், ஓர் இணையுறுப்பு பொய்யாக இருந்தால், மற்றது
வரையறுக்கப்படாததாக இருந்தாலும் இணையல் பொய்யாகும். இதேபோல், பிரிப்பிணைவை
இணைநிலையில் கணிக்கும்போது, இரண்டு பிரிப்புறுப்புகளில் ஒன்றுக்கான செயல்முறை
மெய்யைத் திருப்பித் தந்தவுடன் நிறுத்தலாம்: அப்போது பிரிப்பிணைவு மெய்யாகவே
இருக்க வேண்டும். எனவே ஒரு கூறு வரையறுக்கப்படாததாக இருந்தாலும் கூட்டுக்
கூற்றின் விளைவை இந்நேர்வில் அறியலாம். இப்பொருள்கோளில்~$\Undef$ ஐ
``வரையறுக்கப்படாதது'' என்பதற்குப் பதிலாக ``அறியப்படாதது'' என்று வாசிக்கலாம்.

இவ்விரு பொருள்கோள்கள் க்ளீனியின் வலு தருக்கத்தையும் மெல் தருக்கத்தையும்
தருகின்றன. நிபந்தனைவாக்கியம்~$\lnot !A \lor !B$ க்குச் சமானமாக
வரையறுக்கப்படுகிறது.

\begin{defn}
\emph{வலு க்ளீனி தருக்கம்}~$\LogKs$ பின்வரும் அணியைக் கொண்டு
வரையறுக்கப்படுகிறது:
\begin{enumerate}
  \item $\lnot$, $\land$, $\lor$, $\lif$ ஆகியவற்றைக் கொண்ட செந்தரக்
  கூற்று மொழி~$\Lang L_0$.
  \item மெய்மதிப்புகளின் கணம்~$V = \{\True, \Undef, \False\}$.
  \item $\True$ மட்டுமே நியமிக்கப்பட்ட மதிப்பு; அதாவது,
  $V^+ = \{\True\}$.
  \item வாய்மைச் சார்புகள் பின்வரும் அட்டவணைகளால் தரப்படுகின்றன:
  \begin{center}
    \begin{tabular}{c|c}
      $\tf{\lnot}$ & \\
      \hline
      $\True$ & $\False$ \\
      $\Undef$ & $\Undef$ \\
      $\False$ & $\True$
    \end{tabular}
    \quad
    \begin{tabular}{c|ccc}
      $\tf{\land}[\LogKs]$ & $\True$ & $\Undef$ & $\False$ \\
      \hline
      $\True$ & $\True$ & $\Undef$ & $\False$ \\
      $\Undef$ & $\Undef$ & $\Undef$ & $\False$\\
      $\False$ & $\False$ & $\False$ & $\False$
    \end{tabular}
    \\[2ex]
    \begin{tabular}{c|ccc}
      $\tf{\lor}[\LogKs]$ & $\True$ & $\Undef$ & $\False$ \\
      \hline
      $\True$ & $\True$ & $\True$ & $\True$ \\
      $\Undef$ & $\True$ & $\Undef$ & $\Undef$ \\
      $\False$ & $\True$ & $\Undef$ & $\False$
    \end{tabular}
    \quad
    \begin{tabular}{c|ccc}
      $\tf{\lif}[\LogKs]$ & $\True$ & $\Undef$ & $\False$ \\
      \hline
      $\True$ & $\True$ & $\Undef$ & $\False$ \\
      $\Undef$ & $\True$ & $\Undef$ & $\Undef$  \\
      $\False$ & $\True$ & $\True$ & $\True$
    \end{tabular}
  \end{center}
\end{enumerate}
\end{defn}

\begin{defn}
\emph{மெல் க்ளீனி தருக்கம்}~$\LogKw$ பின்வரும் அணியைக் கொண்டு
வரையறுக்கப்படுகிறது:
\begin{enumerate}
  \item $\lnot$, $\land$, $\lor$, $\lif$ ஆகியவற்றைக் கொண்ட செந்தரக்
  கூற்று மொழி~$\Lang L_0$.
  \item மெய்மதிப்புகளின் கணம்~$V = \{\True, \Undef, \False\}$.
  \item $\True$ மட்டுமே நியமிக்கப்பட்ட மதிப்பு; அதாவது,
  $V^+ = \{\True\}$.
  \item வாய்மைச் சார்புகள் பின்வரும் அட்டவணைகளால் தரப்படுகின்றன:
  \begin{center}
    \begin{tabular}{c|c}
      $\tf{\lnot}$ & \\
      \hline
      $\True$ & $\False$ \\
      $\Undef$ & $\Undef$ \\
      $\False$ & $\True$
    \end{tabular}
    \quad
    \begin{tabular}{c|ccc}
      $\tf{\land}[\LogKw]$ & $\True$ & $\Undef$ & $\False$ \\
      \hline
      $\True$ & $\True$ & $\Undef$ & $\False$ \\
      $\Undef$ & $\Undef$ & $\Undef$ & $\Undef$\\
      $\False$ & $\False$ & $\Undef$ & $\False$
    \end{tabular}
    \\[2ex]
    \begin{tabular}{c|ccc}
      $\tf{\lor}[\LogKw]$ & $\True$ & $\Undef$ & $\False$ \\
      \hline
      $\True$ & $\True$ & $\Undef$ & $\True$ \\
      $\Undef$ & $\Undef$ & $\Undef$ & $\Undef$ \\
      $\False$ & $\True$ & $\Undef$ & $\False$
    \end{tabular}
    \quad
    \begin{tabular}{c|ccc}
      $\tf{\lif}[\LogKw]$ & $\True$ & $\Undef$ & $\False$ \\
      \hline
      $\True$ & $\True$ & $\Undef$ & $\False$ \\
      $\Undef$ & $\Undef$ & $\Undef$ & $\Undef$  \\
      $\False$ & $\True$ & $\Undef$ & $\True$
    \end{tabular}
  \end{center}
\end{enumerate}
\end{defn}

\begin{prop}
  $\LogKs$, $\LogKw$ ஆகிய இரண்டிலும் மெய்மங்கள் இல்லை.
\end{prop}

\begin{proof}
  எல்லா !!{propositional variable}s~$p$ க்கும்
  $\pAssign v(p) = \Undef$ எனில், எந்த வாய்பாடு~$!A$ க்கும்
  மெய்மதிப்பு~$\pValue v(!A) = \Undef$; ஏனெனில் இரண்டு தருக்கங்களிலும்
  \[
    \tf{\lnot}(\Undef) = \tf{\lor}(\Undef, \Undef) = \tf{\land}(\Undef,
  \Undef) = \tf{\lif}(\Undef, \Undef) = \Undef
  \]
  $\LogKs$, $\LogKw$ இரண்டிலும் $\Undef \notin V^+$ என்பதால், இந்த
  !!{valuation} இல்~$!A$ நியமிக்கப்பட்ட மதிப்பைப் பெறாது.
\end{proof}

மெல் க்ளீனி தருக்கத்திலும் வலு க்ளீனி தருக்கத்திலும் மெய்மங்கள் எதுவும்
இல்லை என்றாலும், அவற்றுக்கு அற்பமல்லாத பின்விளைவு உறவுகள் உள்ளன.

\begin{prob}
  பின்வரும் தொடர்புகளில் எவை (a) வலு க்ளீனி தருக்கத்திலும்
  (b) மெல் க்ளீனி தருக்கத்திலும் பொருந்துகின்றன? ஒவ்வொன்றுக்கும் ஒரு மெய்மை
  அட்டவணை தருக.
  \begin{enumerate}
    \item $p, p \lif q \Entails q$
    \item $p \lor q, \lnot p \Entails q$
    \item $p \land q \Entails p$
    \item $p \Entails p \land p$
    \item $p \Entails p \lor q$
  \end{enumerate}
\end{prob}

டிமிட்ரி போச்வார்~$\Undef$ ஐ ``பொருளற்றது'' என்று பொருள்கொண்டார்;
பொய்யுரைப்பவன் முரணுரை போன்ற முரணுரைகளைத் தீர்ப்பதற்காக, முரணுரை
வாக்கியங்கள்~$\Undef$ மதிப்பைப் பெறுகின்றன என்று விதித்து அதைப் பயன்படுத்த
முயன்றார். அடிப்படையில் மெல் க்ளீனி தருக்கத்தைப் போன்ற, ஆனால் கூடுதல்
இணைப்பிகளால் விரிவாக்கப்பட்ட ஒரு தருக்கத்தை அவர் அறிமுகப்படுத்தினார்.
அவற்றில் இரண்டு ``புற மறுப்பு'', ``வரையறுக்கப்படாதது'' செயலி என்பவை:
\begin{center}
  \begin{tabular}{c|c}
    $\tf{\sim}$ & \\
    \hline
    $\True$ & $\False$ \\
    $\Undef$ & $\True$ \\
    $\False$ & $\True$
  \end{tabular}
\quad
  \begin{tabular}{c|c}
    $\tf{+}$ & \\
    \hline
    $\True$ & $\False$ \\
    $\Undef$ & $\True$ \\
    $\False$ & $\False$
  \end{tabular}
\end{center}

\begin{prob}
  போச்வாரின் தருக்கத்தில் $\sim$ ஐ $\lnot$, $+$ ஆகியவற்றைக் கொண்டு
  வரையறுக்க முடியுமா? அதாவது, !!{propositional variable}~$p$ ஐ மட்டுமே
  கொண்டு, $\sim$ இல்லாமல், எப்போதும்~$\mathord{\sim}p$ இன் அதே
  மெய்மதிப்பைப் பெறும் !!a{formula} ஒன்றைக் காண்க. உங்கள் விடை சரி என்பதைக்
  காட்டும் மெய்மை அட்டவணையைத் தருக.
\end{prob}


\end{document}
